\documentclass[11pt,a4paper]{article}
\usepackage[margin=24mm]{geometry}
\usepackage{amsmath,amssymb,amsthm,booktabs}
\newtheorem{lemma}{Lemma}
\newtheorem{theorem}[lemma]{Theorem}
\newcommand{\C}{\mathrm C}\newcommand{\N}{\mathrm N}\newcommand{\Z}{\mathrm Z}
\title{The contribution of projective subgroups to the remaining two counts}
\author{Human Compression project}\date{}
\begin{document}\maketitle
Let $G=O'N$, of order $g=460815505920$. A pair of type $(p,6)$ is
an ordered pair $(r,s)$ with $|r|=p$, $|s|=6$ and $|rs|=2$.
Write $x_J(p,6)$ for the number of these pairs in a subgroup $J$.
Let $\mathcal A$ be their union over the involution centralizers
of $G$, and $\mathcal M$ their union over the two maximal classes
$M\cong L_3(7):2$. We prove
\[
 |\mathcal M(4,6)\setminus\mathcal A|=68g,\qquad
 |\mathcal M(6,6)\setminus\mathcal A|=200g.                 \tag{1}
\]
All subgroup row numbers refer to the acquired complete O'Nan roster.
Its subgroup classification and direct incidence data, the named
character tables, and previously identified projective and affine
models remain explicit borrowed premises. Boolean containment paths
below prove existence only. Every multiplicity we use comes from a
direct incidence or the actual-subgroup arguments given here.

\section*{The internal centralizer union}
In a fixed $M$ there are $2793$ inner involutions, with centralizer
$H$ of order $1344$, and $5586$ outer involutions, with centralizer
$K\cong\mathrm{PGL}_2(7)\times C_2$ of order $672$.
These are the accepted block and orthogonal models, of rows $65/66$
and $110/111$. Fix an inner involution $t$ and put $H=\C_M(t)$.
For $p=4,6$ let $T_i(p)$ count triples $(r,s,u)$ with $(r,s)$ of
the required type in $H$, and $u\ne t$ an inner involution of $H$
centralizing both entries. Define $T_o(p)$ using outer involutions
$u$ instead. The complete literal calculations are
\[
\begin{array}{c|rrrr}
p&x_H(p,6)&x_K(p,6)&T_i(p)&T_o(p)\\\hline
4&1344&2016&0&0\\
6&8736&4704&672&672.
\end{array}                                                \tag{2}
\]
The reader reconstructs each finite group from its retained generators,
partitions the order-six elements into complete conjugacy classes,
and every partner set into orbits under the full centralizer of its
representative. Orders, common involutions and generated subgroups
are checked by literal permutation composition. Inner versus outer
is the preserved point/line partition in the 114-point model.

\begin{lemma}\label{internal}
The number $U_M(p)$ of pairs centralized by an involution inside $M$ is
\[
 U_M(p)=2793\bigl(x_H(p,6)-T_i(p)/3-T_o(p)\bigr)
                      +5586x_K(p,6).                       \tag{3}
\]
Its values for $p=4,6$ are $15015168$ and $48173664$.
\end{lemma}
\begin{proof}
The accepted centralizer lemma gives
$\C_G(s)\cong C_3^2\times D_8$. Its Sylow 2-subgroup is normal
and has elementary abelian rank at most two. The same is therefore true
of the common centralizer inside $M$. The classical elementary
abelian subgroup cancellation identity counts its involutions
with coefficient $1$ and its four-groups with coefficient $-2$.
The index-two socle of $M$ meets each four-group in either the whole
four-group or one involution.

A four-group with three inner involutions is counted six times
in the sum of $T_i$ over inner $t$: three choices of $t$ and two
choices of $u$. Its required correction $-2$ thus has coefficient
$-1/3$. A four-group with just one inner involution is counted twice
in the sum of $T_o$, giving coefficient $-1$. Every four-group has
one of these two forms. Summing over the two involution classes
proves (3), and (2) gives the stated integers.
\end{proof}

\begin{lemma}\label{external}
A pair of either type in $M$ belongs to $\mathcal A$ if and only
if it is centralized by an involution inside $M$.
\end{lemma}
\begin{proof}
Only the forward implication needs proof. Put $L=\langle r,s\rangle$.
The complete roster shows that every possible class of $L$ below
both row $7$ and row $2$ or $3$ is below one of $65,66,110,111$.
Conjugate $L$ into such a block or orthogonal subgroup. The literal
partition used in (2) additionally closes each generated subgroup
and supplies a central involution in it, except for the centerless
generated subgroups of order $336$. A central involution of $L$
conjugates back into $L\le M$, proving the assertion in all those
cases. This use requires no alignment of the chosen block with $M$.

The possible order-$336$ classes here are $171,173,175,176,177$.
The latter three have a central involution, by their accepted
$\mathrm{SL}_2(7)$ or $\mathrm{PSL}_2(7)\times C_2$ structures.
For class $171$ or $173$, reuse the established actual-parent lemma:
$L$ has a unique parent $K$ in row $110$ or $111$ and a unique
simple parent $S$ in row $4$ or $5$. Since $L\lhd K$ and
$|\N_G(L)|=672$, $K=\N_G(L)$ normalizes $S$. Therefore
$K\le\N_G(S)$, the unique projective maximal containing $S$.
The complete parent lists show that every projective overgroup
of $L$ is this maximal and hence contains $K$. Its central
involution centralizes $L$ inside $M$. We need existence here,
not uniqueness of the global common involution.
\end{proof}

\section*{Aligning the small actual overgroups}
Each projective maximal has two maximal subgroup types of order $144$:
$N\cong M_9:2$, with normal Sylow subgroup $C_3^2$, and
$D\cong S_3\times S_4$. They are rows $250/252$ and $251/253$.
For clarity, the model of $N$ is the retained order-$144$ subgroup
generated by a type-$(4,8)$ pair in the already identified affine
group $C$. Its roster join places it in $250/252$. The two copies
of $L_3(7):2$ have the same abstract maximal types: $N$ has elements
of order eight and $D$ does not, distinguishing the two order-$144$
types in both classes. The accepted structural argument identifies
$D$ as $S_3\times S_4$, also reconstructed in its natural disjoint
action. A fresh literal calculation on these two small models gives
\[
\begin{array}{c|cc}
&p=4&p=6\\\hline
N&72\ (144)&36\ (144)\\
D&\text{none}&36\ (288).
\end{array}                                                \tag{4}
\]
Only pairs without an intrinsic central involution are listed;
the entries are generated orders, followed by counts. Every generated
group in (4) has a normal Sylow subgroup of order nine, checked as
the closed set of its identity and all order-three elements.

\begin{lemma}\label{small}
Every pair of either type outside $\mathcal A$ whose generated
subgroup lies in one of the four order-$144$ classes has exactly
one projective maximal overgroup.
\end{lemma}
\begin{proof}
By (4), its generated subgroup $L$ has order $72$ or $36$, with
characteristic Sylow subgroup $V\cong C_3^2$. In either projective
class the only subgroup class of order nine is $553$ or $554$,
with global normalizer of order $1296$. The retained $N$ normalizes
a Sylow subgroup of order nine, so Sylow conjugacy and
$\gcd(1296,3753792)=144$ show that $\N_M(V)$ is an actual copy
of the appropriate $N$, of order $144$. As $V\lhd L$, every
projective maximal containing $L$ also contains such an $N$
containing $L$.

The complete roster interval now has only the following possibilities.
For order $72$, $L$ is row $330$ or $331$ and has one actual parent
in $252$ or $250$. For order $36$, $L$ is row $406$ or $407$;
the intervals to $252$ and $250$ are respectively
$406<330<252$ and $407<331<250$. Each edge gives one actual parent:
the parent class length times its displayed child multiplicity
equals the child class length. Finally each $250/252$ subgroup
has one actual projective maximal parent. Thus all possible
projective maximal overgroups coincide. In particular the multiple
$S_3\times S_4$ routes above $S_3\times S_3$ cannot be added as
independent overgroups.
\end{proof}

\section*{The flag normalizer and repeated projective parents}
Let $B$ be row $22$, the flag normalizer of order $8232$.
Its normal Sylow 7-subgroup $U$ has order $343$ and is row $170$;
$B=\N_G(U)$. The literal retained flag model gives
\[
x_B(4,6)=0,\qquad x_B(6,6)=403368,
\qquad |B(6,6)\setminus\mathcal A|=296352=36|B|.            \tag{5}
\]
Indeed its internally covered pairs number $107016$, generating
orders $12$ or $84$. Its remaining pairs generate orders $588$
or $4116$, with counts $197568$ and $98784$. The only corresponding
classes below $B$ are $134,135$ and $27,28,29$; none is below row $7$.
Consequently this internal partition is exactly the global
$\mathcal A$ partition asserted in (5).

\begin{lemma}\label{flag}
If $L$ is generated by one of the uncovered pairs in (5), then
the number of its actual projective maximal overgroups minus one
equals its number of actual flag-normalizer overgroups.
For $|L|=4116$ these numbers of overgroups are $2$ and $1$;
for $|L|=588$ they are $3$ and $2$.
\end{lemma}
\begin{proof}
Suppose first that $|L|=4116$. It has normal Sylow subgroup $U$
of order $343$. This is also Sylow in every projective maximal
$M$ containing $L$. The projective Sylow normalizer has order
$8232$, equal to the full global normalizer of row $170$.
Thus every such $M$ contains the same $B=\N_G(U)$.
Conversely $B$ has exactly one actual parent in each projective
class, by its direct incidences. There is just one flag normalizer
containing $L$, since its Sylow subgroup must be $U$. This proves
the assertion for rows $27,28,29$.

Now take row $134$; row $135$ is symmetric. Its Sylow subgroup
$V$ of order $49$ is unique, since its number divides $12$ and
is $1$ modulo $7$. There are exactly two actual row-$27$ parents
$D_1,D_2$ and one actual row-$82$ parent $N$, of order $1176$.
As $L$ has index two in $N$, $N$ normalizes $L$ and its $V$.
The three order-$49$ classes have global normalizer orders
$32928,2058,32928$. Divisibility by $1176$ excludes $2058$.
The only subgroup classes of order $32928$ are the parabolics
in rows $9,10$. Since $N$ has a unique actual parabolic parent
$P$, this gives $P=\N_G(V)$.

For each $D_i$, let $U_i$ be its normal Sylow subgroup of order
$343$. It contains $V$, which has index seven in the 7-group
$U_i$ and is therefore normal in $U_i$. Since $L$ normalizes $V$
and $D_i=U_iL$, also $D_i$ normalizes $V$. Both $D_i$ therefore
lie in $P$, whose unique projective parent $M_0$ is of row $2$.
Each $D_i$ has its unique flag parent $B_i$. The two $B_i$ are
distinct, because a flag normalizer contains just one row-$27$
subgroup. By the order-$4116$ argument, both $B_i$ lie in $M_0$.
Their opposite projective parents, of row $3$, are distinct:
otherwise two distinct flag normalizers would lie in both $M_0$
and that opposite maximal. A flag normalizer is maximal in $M_0$,
so the two generate $M_0$, an impossibility.
The complete immediate-parent list of $L$ is just $27,82$;
the route through $82$ leads to the same $M_0$. Thus there are
exactly three projective parents and two flag parents.
\end{proof}

\begin{lemma}\label{exhaust}
Every remaining pair outside $\mathcal A$ in a projective maximal,
whose generated subgroup is not one of the classes in
Lemma~\ref{flag}, has a unique projective maximal overgroup.
\end{lemma}
\begin{proof}
Intersect the complete ancestors of the cyclic classes of the two
orders with the projective descendants, and remove row-$7$
descendants. Separate flag descendants and descendants of the
four order-$144$ classes. For $p=4$ the remaining list is
\[
2,3,4,5,9,10,13,14,45,46,82,83,84,85,136,137,172,174.        \tag{6}
\]
For $p=6$, add $138,139,285$. Row $285$ is $19:6$ and has no
pair of type $(6,6)$: both factors inject as generators of its
cyclic quotient of order six, whose product has odd order;
an involution injects through the odd kernel. Flag descendants
are fully settled by (5), and the small descendants by
Lemma~\ref{small}.

In (6), all classes except $136,137$ lead by direct unique-parent
edges to one projective maximal. In particular $138,139$ have
unique parents $82,83$. A row-$136$ subgroup $L$ has order $588$
and unique Sylow subgroup $V$ of order $49$. Any parent of class
$45$ has a normal Sylow subgroup of order $49$, necessarily $V$;
the same is true of the class-$82$ parent. These are all its
immediate-parent types. The displayed normal $7^2$ kernels in rows
$45,46,82,83$ are borrowed structural roster premises. They normalize
$V$, whose full normalizer
is the unique parabolic of order $32928$, since its order is
divisible by $1176$ and not $2058$. All routes therefore lead
to that parabolic's same projective maximal. Row $137$ is symmetric.
This proves the assertion for every class in the complete lists.
\end{proof}

\begin{theorem}
The projective union has the values in (1). Together with the
accepted complementary-family proof it gives the two source
entries of types $(4,6)$ and $(6,6)$ by exact character subtraction.
\end{theorem}
\begin{proof}
Each of the two projective classes has $122760$ members. The
acquired character values for one such maximal are
\[
x_M(4,6)=142644096,\qquad x_M(6,6)=491121120.
\]
By Lemmas~\ref{internal} and \ref{external}, subtracting $U_M$
counts exactly its pairs outside $\mathcal A$. Lemmas~\ref{flag}
and \ref{exhaust} prove that the excess in the sum over all
projective maximals is exactly the sum over flag normalizers of
their uncovered pairs. Therefore
\[
|\mathcal M(p,6)\setminus\mathcal A|
=2\cdot122760\bigl(x_M(p,6)-U_M(p)\bigr)
-\frac{g}{8232}\begin{cases}0,&p=4,\\296352,&p=6.\end{cases}
\]
The first term is $68g$ or $236g$, and the last is zero or $36g$.
This proves (1).

The separately accepted family-coverage proof gives $2g$ and
$46g$ outside $\mathcal A\cup\mathcal M$. The accepted common
centralizer formula gives $|\mathcal A|=17g/2$ and $113g/9$.
Subtracting these disjoint nongenerating contributions from the
acquired full-group character counts, then dividing by
$|\operatorname{Aut}(G)|=2g$, gives the ordered automorphism-orbit
entries
\[
 \frac{75804150723840-(17/2+68+2)g}{2g}=43,\qquad
 \frac{271215526095360-(113/9+200+46)g}{2g}=165.
\]
The action is free on generating pairs; no reversal
quotient or chirality conclusion is included.
\end{proof}

The accompanying certificate checks every stated roster set,
normalizer, direct incidence, literal partition and rational
identity. The source numerical values are attributed to Connor
and Leemans; the present proof uses their subgroup classification
as a premise and supplies an independently replayable assembly.
This component gives no whole-original, new conjecture or Lean credit.
\end{document}
